\section{Conclusion}\label{conclusion}

Across approximately 21,540 runs spanning 12 models from 11 independent training pipelines, the Conceptual Topology Mapping Benchmark establishes that LLMs have consistent, measurable geometric structure in how they navigate between concepts. This structure is consistent with quasimetric geometry --- systematically asymmetric (mean 0.811), with the triangle inequality holding at approximately 91\%. Navigation is organized around bottleneck bridges that function as structural landmarks, not mere associations, and expressed through model-specific gaits spanning a 2.9x range from GPT (0.258) to Mistral (0.747). These findings generalize across all 12 models and resist protocol variation.

The mechanism ceiling is itself the most important finding. Seven follow-up hypotheses --- spanning mechanistic, model-deficit, methodological, and effect-direction types --- were falsified across Phases 8--10; one was resurrected with additional data. Prediction accuracy descended from 81.3\% for characterization questions to 20--24\% for mechanistic ones --- not because signal disappeared, but because qualitative structure (characterizable at \textasciitilde80\%) resists quantitative mechanistic decomposition (opaque at 0--15\%). The dynamics generating navigational regularity do not reduce to any individual variable tested. This parallels the broader challenge in interpretability: structure is detectable before it is explainable.

Static embedding geometry is well-studied; navigation on that geometry is not. \citet{karkada2025translation} proved why the geometric structure exists --- translation symmetry in co-occurrence statistics determines manifold shape. We showed what that geometry supports: structured, asymmetric, model-specific behavioral navigation. The Conceptual Topology Mapping Benchmark provides a different evaluation paradigm --- not testing what models know, but how they navigate what they know. The 29 documented dead ends and the mechanism ceiling are not failures of the research program but honest accounting of where the field's explanatory tools fall short, and where the next generation of multi-factorial mechanistic models will need to reach.
